\documentclass[
	10pt,								% globale Schriftgröße
	parskip=half-,						% setzt Absatzabstand hoch
	paper=a4,							% Format
	english,ngerman,					% lädt Sprachpakete
	]{scrartcl}							% Dokumentenklasse

% //////////////////// Pakete laden ////////////////////
\usepackage[fleqn]{amsmath}
\usepackage[fleqn]{mathtools}
\usepackage{amssymb}			% mathematische symbole, für \ceckmarks
\usepackage{amsthm}				% für proof
\usepackage{mathrsfs}			% für \mathscr
\usepackage{latexsym}
\usepackage{marvosym}				% für Lightning

\usepackage{fontspec} 			% funktioniert nur mit den neueren Compilern z.B. XeLaTeX
\usepackage{microtype}			% für bessere Worttrennung
\usepackage[ngerman]{babel} 	% Spracheinstellung
\usepackage{lmodern}			% verändert verwendete Schriftart, damit sie weniger pixelig ist

\usepackage{verbatim}
\usepackage{listings}			% Für Quellcode

\usepackage{graphicx}
\usepackage{tabularx}			% für Tabellen mit gleicher Spaltenbreite und automatischen Umbrüchen
\usepackage{fullpage}
\usepackage{multirow}			% für multirow in tabulars
\usepackage{rotate}
\usepackage[cmyk,table]{xcolor} % um Farben zu benutzen, kann mehr als das Paket color
\usepackage[					% Verlinkungen
	colorlinks,					% farbige Schrift, statt farbiger Rahmen
	linktocpage,				% verlinkt im Abb.Verzeichnis Seitenzahl statt Bildunterschrift
	linkcolor=blue				% setzt Farbe der Links auf blau
	]{hyperref}					% nur für digitale Anwendungen, url = "http://www.example.com"
\usepackage{url}				% für Webadressen wie e-mail usw.: "\url{http://www.example.com}"

\usepackage{enumerate}			% für versch. Aufzählungezeichen wie z.B. a)
\usepackage{xspace}				% folgt ein Leerzeichen nach einem \Befehl, wird es nicht verschluckt.
\usepackage{cancel}				% für das Durchstreichen u.a. in Matheformeln mit \cancel
\usepackage{float}              % zum Forcieren der Position von figure-Umgebungen

% zum Zeichnen (u.a. von Graphen)
\usepackage{fp}
\usepackage{tikz}
\usetikzlibrary{tikzmark}			% für \tikzmark{toRemember}
\usetikzlibrary{positioning}	% verbesserte Positionierung der Knoten
\usetikzlibrary{automata}		% für Automaten (GTI)
\usetikzlibrary{arrows}
\usetikzlibrary{shapes}
\usetikzlibrary{decorations.pathmorphing}
\usetikzlibrary{decorations.pathreplacing}
\usetikzlibrary{decorations.shapes}
\usetikzlibrary{decorations.text}

% //////////////////// Syntaxhighlighting ////////////////////
\lstloadlanguages{Python, Haskell, [LaTeX]TeX, Java}
\lstset{
   basicstyle=\footnotesize\ttfamily,	% \scriptsize the size of the fonts that are used for the code
   backgroundcolor = \color{bgcolour},	% legt Farbe der Box fest
   breakatwhitespace=false,	% sets if automatic breaks should only happen at whitespace
   breaklines=true,			% sets automatic line breaking
   captionpos=t,				% sets the caption-position to bottom, t for top
   commentstyle=\color{codeblue}\ttfamily,% comment style
   frame=single,				% adds a frame around the code
   keepspaces=true,			% keeps spaces in text, useful for keeping indentation
							% of code (possibly needs columns=flexible)
   keywordstyle=\bfseries\ttfamily\color{codepurple},% keyword style
   numbers=left,				% where to put the line-numbers;
   							% possible values are (none, left, right)
   numberstyle=\tiny\color{codegreen},	% the style that is used for the line-numbers
   numbersep=5pt,			% how far the line-numbers are from the code
   stepnumber=1,				% nummeriert nur jede i-te Zeile
   showspaces=false,			% show spaces everywhere adding particular underscores;
							% it overrides 'showstringspaces'
   showstringspaces=false,	% underline spaces within strings only
   showtabs=false,			% show tabs within strings adding particular underscores
   flexiblecolumns=false,
   tabsize=1,				% the step between two line-numbers. If 1: each line will be numbered
   stringstyle=\color{orange}\ttfamily,	% string literal style
   numberblanklines=false,				% leere Zeilen werden nicht mitnummeriert
   xleftmargin=1.2em,					% Abstand zum linken Layoutrand
   xrightmargin=0.4em,					% Abstand zum rechten Layoutrand
   aboveskip=2ex, 
}

\lstdefinestyle{py}{
   language=Python,
}
\lstdefinestyle{hs}{
   language=Haskell,
}
\lstdefinestyle{tex}{
	language=[LaTeX]TeX,
	escapeinside={\%*}{*)},     % if you want to add LaTeX within your code
	texcsstyle=*\bfseries\color{blue},% hervorhebung der tex-Schlüsselwörter
	morekeywords={*,$,\{,\},\[,\],lstinputlisting,includegraphics,
	rowcolor,columncolor,listoffigures,lstlistoflistings,
	subsection,subsubsection,textcolor,tableofcontents,colorbox,
	fcolorbox,definecolor,cellcolor,url,linktocpage,subtitle,
	subject,maketitle,usetikzlibrary,node,path,addbibresource,
	printbibliography},% if you want to add more keywords to the set
     numbers=none,
     numbersep=0pt,
     xleftmargin=0.4em,
}

\lstdefinestyle{java}{
	language=Java,
	extendedchars=true,		% lets you use non-ASCII characters;
   						% for 8-bits encodings only, does not work with UTF-8
}

\lstdefinelanguage[x64]{Assembler}     % add a "x64" dialect of Assembler
   [x86masm]{Assembler} % based on the "x86masm" dialect
   % with these extra keywords:
   {morekeywords={CDQE,CQO,CMPSQ,CMPXCHG16B,JRCXZ,LODSQ,MOVSXD, %
                  POPFQ,PUSHFQ,SCASQ,STOSQ,IRETQ,RDTSCP,SWAPGS, %
                  rax,rdx,rcx,rbx,rsi,rdi,rsp,rbp, %
                  r8,r8d,r8w,r8b,r9,r9d,r9w,r9b}
}					% for 8-bits encodings only, does not work with UTF-8

\lstdefinestyle{c}{
	language=c,
	extendedchars=true,		% for 8-bits encodings only, does not work with UTF-8
}

% //////////////////// eigene Kommandos ////////////////////
\newcommand\FU{Freie Universität Berlin\xspace}% benötigt package xspace
\newcommand\gdw{g.\,d.\,w.\xspace}
\newcommand\oBdA{o.\,B.\,d.\,A.\xspace}
\newcommand{\Eu}{\texteuro}
\newcommand\N{\mathbb{N}\xspace}
\newcommand\Q{\mathbb{Q}\xspace}
\newcommand\R{\mathbb{R}\xspace}
\newcommand\Z{\mathbb{Z}\xspace}
\newcommand\ohneNull{\ensuremath{\backslash\lbrace 0\rbrace}}% \{0}
\let\dhALT\dh	% Schreibt Befehl \dh in \dhALT um
\renewcommand\dh{d.\,h.\xspace}	%renew überschreibt command \dh
\newcommand\Bolt{\;\text{\LARGE\raisebox{-0.3em}{\Lightning}\normalsize}\xspace}% Blitz
\newcommand\zz{\ensuremath{\raisebox{+0.25ex}{Z}% zu zeigen
			\kern-0.4em\raisebox{-0.25ex}{Z}%
			\;\xspace}}
\newcommand{\from}{\ensuremath{\colon}}
\newcommand{\floor}[1]{\lfloor{#1}\rfloor}
\newcommand{\ceil}[1]{\lceil{#1}\rceil}
 \renewcommand{\L}{\ensuremath{\mathcal{L}}\xspace}
 \renewcommand{\P}{\ensuremath{\mathcal{P}}\xspace}
 \newcommand{\NL}{\ensuremath{\mathcal{N}\kern-0.2em\mathcal{L}}\xspace}
 \newcommand{\NP}{\ensuremath{\mathcal{NP}}\xspace}

% //////////////////// Mathefunktionen ////////////////////
\DeclareMathOperator{\Landau}{\mathcal{O}}
\DeclareMathOperator{\True}{True}
\DeclareMathOperator{\False}{False}

% //////////////////// eigene Theoreme ////////////////////
\newtheorem{theorem}{Satz}
\newtheorem{corollary}[theorem]{Folgerung}
\newtheorem{lemma}[theorem]{Lemma}
\newtheorem{observation}[theorem]{Beobachtung}
\newtheorem{definition}[theorem]{Definition}
\newtheorem{Literatur}[theorem]{Literatur}
% konfiguriert proof
\makeatletter
\newenvironment{Proof}[1][\proofname]{\par
  \pushQED{\qed}%
  \normalfont \topsep6\p@\@plus6\p@\relax
  \trivlist
  \item[\hskip\labelsep
%         \itshape
        \bfseries
    #1\@addpunct{.}]\ignorespaces
}{%
  \popQED\endtrivlist\@endpefalse
}
\makeatother

% //////////////////// eigene Farben ////////////////////
\let\definecolor=\xdefinecolor
\definecolor{FUgreen}{RGB}{153,204,0}
\definecolor{FUblue}{RGB}{0,51,102}

\definecolor{middlegray}{rgb}{0.5,0.5,0.5}
\definecolor{lightgray}{rgb}{0.8,0.8,0.8}
\definecolor{orange}{rgb}{0.8,0.3,0.3}
\definecolor{azur}{rgb}{0,0.7,1}
\definecolor{yac}{rgb}{0.6,0.6,0.1}
\definecolor{Pink}{rgb}{1,0,0.6}

\definecolor{bgcolour}{rgb}{0.97,0.97,0.97}
\definecolor{codegreen}{rgb}{0,0.6,0}
\definecolor{codegray}{rgb}{0.35,0.35,0.35}
\definecolor{codepurple}{rgb}{0.58,0,0.82}
\definecolor{codeblue}{rgb}{0.4,0.5,1}

% //////////////////// eigene Settings ////////////////////

\textheight = 230mm		% Höhe des Satzspiegels / Layouts
\footskip = 10ex			% Abstand zw. Fußzeile und Grundlinie letzter Textzeile
\parindent 0pt			% verhindert Einrückung der 1. Zeile eines Absatzes
\setkomafont{sectioning}{\rmfamily\bfseries}% setzt Ü-Schriften in Serifen, {disposition}											% bindet Header ein (WICHTIG)
\usepackage{graphicx}
\usepackage{amsmath}
\usepackage{amssymb}
\usepackage{fancyvrb}

\newcommand{\dozent}{Prof. R. Rojas}					% <-- Names des Dozenten eintragen
\newcommand{\projectNo}{2}
\newcommand{\veranstaltung}{Mustererkennung}
\newcommand{\semester}{WS17/18}
\newcommand{\studenten}{Boyan Hristov, Nedeltscho Petrov}
% /////////////////////// BEGIN DOKUMENT /////////////////////////


\begin{document}
% /////////////////////// BEGIN TITLEPAGE /////////////////////////
\begin{titlepage}
	\subject{\dozent}
	\title{\veranstaltung, \semester}
	\subtitle{\Large Übungsblatt \projectNo\\ \large\vspace{1ex} TutorIn: \tutor\\ Tutorium \tutoriumNo}
	\author{\studenten}
	\date{\normalsize \today}
\end{titlepage}

\maketitle								% Erstellt das Titelblatt
\vspace*{-10cm}							% rückt Logo an den oberen Seitenrand
\makebox[\dimexpr\textwidth+1cm][r]{	%rechtsbündig und geht rechts 1cm über Layout hinaus
	\includegraphics[width=0.4\textwidth]{src/fu_logo} % fügt FU-Logo ein
}
% /////////////////////// END TITLEPAGE /////////////////////////

\vspace{7cm}							% Abstand
\rule{\linewidth}{0.8pt}				% horizontale Linie										% erstellt die Titelseite


Link zum Git Repository: \url{https://github.com/BoyanH/Freie-Universitaet-Berlin/tree/master/MachineLearning/Homework\projectNo}

% /////////////////////// Aufgabe 1 /////////////////////////

\section*{Score}
Das ist die Ausgabe des Programs und damit auch das Score
\begin{lstlisting}
Score for 3 vs 7: 100.0%
Score for 3 vs 8: 100.0%
Score for 5 vs 7: 100.0%
Score for 5 vs 8: 100.0%
Score for 7 vs 8: 100.0%
\end{lstlisting}

\section*{Lineare Regression}

Wir benutzen die folgende Formel für lineare Regression mit Least Squares
\begin{align}
	\vec{B} = (X^TX)^{-1}X^T\vec{y}
\end{align}
wobei $\vec{B}$ die Lösung von der folgenden Gleichung ist: (im besten Fall, falls Matrix mit vollem Rank)

\begin{align}
	y_i = B_0 + B_iX_{11} + B_iX{12} + ... + B_nX_{in}
\end{align}

dabei ist X die Eingabematrix mit je Zeile ein Punkt im n-Dimensionalen Raum.
Damit wir diese Gleichung lösen können, mit $B_0$ als 1. Glied, müssen wir erstmal eine Spalte mit Einsen einfügen.

\begin{lstlisting}[style=py]
	ones = np.ones((len(self.trainData), 1), dtype=float)
	X = np.append(ones, self.trainData, axis = 1)
\end{lstlisting}

Weiter, um die 2 Klassen gleichwertig in der linearen Regression zu betrachten, dürfen wir natürlich nicht gleich die Labels benutzen, sondern
diese erstmal normalisieren (so zu sagen). Dabei mapen wir die eine Labels zu -1 und die andere zu 1. Z.B wenn wir zwischen 3 und 5 unterscheiden wollen,
wird das je Label 3 zu -1 und je 5 zu 1. Das passiert folgendermaßen

\begin{lstlisting}[style=py]
	def normalizeLabels(self, labels):
		return list(map(lambda x: -1 if int(x) == self.classA else 1, labels))
\end{lstlisting}

Und damit unsere fit Methode

\begin{lstlisting}[style=py]
	def fit(self):
		ones = np.ones((len(self.trainData), 1), dtype=float)
		X = np.append(ones, self.trainData, axis = 1)
		xtxInversed = LinearRegressionClassifier.pseudoInverse(X.T.dot(X))
		normalizedLabels = self.normalizeLabels(self.trainLabels)
		self.beta = xtxInversed.dot(X.T).dot(normalizedLabels)
\end{lstlisting}

\section*{Pseudoinverse}
Wir benutzen die Moore-Penrose Pseudoinverse. Nach der Folmel gilt:

\begin{align}
	A^{+} = \lim_{\delta \rightarrow 0} A^{*}(AA^{*} + \delta E)^{-1}
\end{align}

Wobei $A^{*}$ die adjungierte Matrix ist. Da wir aber in dem Bereich der reelen Zahlen uns befinden, ist die adjungierte Matrix
identisch zu der transponierten. Deswegen gilt:

\begin{align}
	A^{+} = \lim_{\delta \rightarrow 0} A^{T}(AA^{T} + \delta E)^{-1}
\end{align}

Damit unsere Implementierung

\begin{lstlisting}[style=py]
@staticmethod
	def pseudoInverse(X):
		delta = np.nextafter(np.float16(0), np.float16(1)) # as close as we can get to lim delta -> 0
		pseudoInverted = X.T.dot(np.linalg.inv(X.dot(X.T) + delta * np.identity(len(X))))
\end{lstlisting}


\section*{Vollständige Implementierung (Ohne Classifier Klasse)}
\begin{lstlisting}[style=py]
import pandas as pd
import numpy as np
from operator import itemgetter
import matplotlib.pyplot as plt
import os
from Classifier import Classifier

import seaborn as sn
import pandas as pd
import matplotlib.pyplot as plt

class LinearRegressionClassifier(Classifier):


	# B = (X^TX)^-1X^T*y

	# hat H = X*B^ =` X(X^TX)-1X^T

	@staticmethod
	def pseudoInverse(X):
		# Determining whether a matrix is invertable or not with gaus elimination
		# (most efficient from the naive approaches) takse n^3 time, which
		# makes the whole programm slower for not much precision gain
		# therefore we simply calculate the pseudo invsersed matrix every time 

		# if LinearRegressionClassifier.isInvertable(X):
		# 	return np.linalg.inv(X)

		# calculate pseudo-inverse A+ of a matrix A (X in our case)
		# A+ = lim delta->0 A*(A.A* + delta.E)^(-1) where A* is the conjugate transpose
		# and E ist the identity matrix

		# In our case, we are working with real numbers, so A* = A^T
		# so the formula is A^T(A.A^T + delta.E)^(-1)

		delta = np.nextafter(np.float16(0), np.float16(1)) # as close as we can get to lim delta -> 0
		pseudoInverted = X.T.dot(np.linalg.inv(X.dot(X.T) + delta * np.identity(len(X))))

		return pseudoInverted

	@staticmethod
	def isInvertable(X):
		# apply gaus elimination
		# if the matrix is transformable in row-echelon form
		# then it is as well inverable

		X = np.copy(X) # don't really change given matrix
		m = len(X)
		n = len(X[0])
		for k in range(min(m, n)):
			# Find the k-th pivot:
			# i_max  = max(i = k ... m, abs(A[i, k]))
			i_max = k
			max_value = X[k][k]
			for i in range(k, m):
				if X[i][k] > max_value:
					max_value = X[i][k]
					i_max = i
			if X[i_max][k] == 0:
				return False
			for i in range(n):
				temp = X[k][i]
				X[k][i] = X[i_max][i]
				X[i_max][i] = temp
			# Do for all rows below pivot:
			for i in range(k+1, m):
				# for i = k + 1 ... m:
				f = X[i][k] / X[k][k]
				# Do for all remaining elements in current row:
				for j in range(k + 1, n):
					X[i][j]  = X[i][j] - (X[k][j] * f)
				X[i][k] = 0

		return True

	def __init__(self, trainSet, testSet, classA, classB):
		self.classA = classA
		self.classB = classB

		trainSet = self.filterDataSet(trainSet)
		testSet = self.filterDataSet(testSet)

		self.trainData = list(map(lambda x: x[1:], trainSet))
		self.trainLabels = list(map(itemgetter(0), trainSet))
		self.testSet = list(map(lambda x: x[1:], testSet))
		self.testLabels = list(map(itemgetter(0), testSet))

		self.fit()

	def filterDataSet(self, dataSet):
		return list(filter(lambda x: int(x[0]) in [self.classA, self.classB], dataSet))

	def fit(self):
		# fill X with (1,1...,1) in it's first column to be able to get the
		# wished yi = B0 + B1Xi1 + B2Xi2 + ... + BnXin
		ones = np.ones((len(self.trainData), 1), dtype=float)
		X = np.append(ones, self.trainData, axis = 1)

		# and then used the following formula to calculate our closest possible B
		# which solves best our least squares regression
		# B = (X^TX)^(-1)X^Ty
		# where y = (-1, 1, -1, 1, ..., -1) (for example) is a vector
		# of the labels corresponding to the given data points

		xtInversed = LinearRegressionClassifier.pseudoInverse(X.T.dot(X))

		# normalize y, so the two possible classes are maped to -1 or 1
		normalizedLabels = self.normalizeLabels(self.trainLabels)
		self.beta = xtInversed.dot(X.T).dot(normalizedLabels)

	def predictSingle(self, X):
		X = np.append(np.array([1]), np.array(X), axis=0)
		return self.classA if (X.dot(self.beta) < 0) else self.classB

	def predict(self, X):
		return np.array(list(map(lambda x: self.predictSingle(x), X)))

	def test(self):
		print('Score for {} vs {}: {}%'.format(
			self.classA, self.classB, self.score(self.testSet, self.testLabels) * 100))
		self.printConfusionsMatrix(self.confusion_matrix(self.testSet, self.testLabels))


	def normalizeLabels(self, labels):
		return list(map(lambda x: -1 if int(x) == self.classA else 1, labels))

	def printConfusionsMatrix(self, matrix):
		explicitImgPath = os.path.join(dir_path, './Plots/confusion_matrix_for_{}vs_{}.png'.format(
			self.classA, self.classB))
		digits = [str(x) for x in range(10)];

		df_cm = pd.DataFrame(matrix, index = digits,
		columns = digits )

		plt.figure(figsize = (11,7))
		heatmap = sn.heatmap(df_cm, annot=True)

		heatmap.set(xlabel='Klassifiziert', ylabel='Erwartet')

		plt.savefig(explicitImgPath, format='png')

def extractDataFromLine(line):
	line = line.replace(' \n', '') # clear final space and new line chars
	return list(map(float, line.split(' '))); # map line to a list of floats

def parseDataFromFile(file, dataArr):
	for line in file:
		line = line.replace(' \n', '') # clear final space and new line chars
		currentDigitData = extractDataFromLine(line)
		dataArr.append(currentDigitData)


trainSet = []
testSet = []

dir_path = os.path.dirname(os.path.realpath(__file__))
explicitPathTrainData = os.path.join(dir_path, './Dataset/train')
explicitPathTestData = os.path.join(dir_path, './Dataset/test')
trainFile = open(explicitPathTestData, 'r')
testFile = open(explicitPathTestData, 'r')

parseDataFromFile(trainFile, trainSet)
parseDataFromFile(testFile, testSet)

LinearRegressionClassifier(trainSet, testSet, 3, 5).test()
LinearRegressionClassifier(trainSet, testSet, 3, 7).test()
LinearRegressionClassifier(trainSet, testSet, 3, 8).test()
LinearRegressionClassifier(trainSet, testSet, 5, 7).test()
LinearRegressionClassifier(trainSet, testSet, 5, 8).test()
LinearRegressionClassifier(trainSet, testSet, 7, 8).test()
\end{lstlisting}



% /////////////////////// END DOKUMENT /////////////////////////
\end{document}